\documentclass{situationist-preprint}

\title{Letters of Recommendation and Quantity Implicature}
\msid{LNG-11-b}
\msfamily{LNG-11}
\msdiscipline{Philosophy of Language}
\mstype{Companion manuscript}
\mskeywords{implicature; Quantity; flouting; recommendation letters; non-lexical alternatives; competence}

\begin{document}
\maketitle

\begin{abstract}
Grice's testimonial for a philosophy candidate, which praises only his English and his attendance, conveys that he is no good at philosophy. We analyse the case as a flouting of Quantity and show that it is derived by the two steps of LNG-11-a: the writer does not believe the stronger alternative, and, being competent, believes its negation. The alternative is not a scale-mate of anything the writer said; it is supplied by the purpose of the exchange. The case shows that the second step applies to non-lexical alternatives, and that whether it applies depends on what the writer knows, which only a pragmatic derivation can track.
\end{abstract}

\section{The case}

\citet[52]{Grice1975} imagines a professor writing a testimonial for a pupil who is a candidate for a philosophy post: \enquote{Mr X's command of English is excellent, and his attendance at tutorials has been regular.} The reader infers that the writer thinks Mr X is no good at philosophy. Grice classes the case as a flouting of the first maxim of Quantity: the writer cannot be opting out, since she wrote the letter at all; she cannot be ignorant, since he is her pupil; so she must wish to convey something she is reluctant to write down.

LNG-11-a derives scalar implicatures in two steps. Proposition 5.1 of LNG-11-a gives the primary implicature: a speaker presumed to observe Quantity, Quality and Manner, who says $\phi$, does not believe any stronger alternative of $\phi$. Proposition 5.2 gives the secondary implicature: if she is competent with respect to the alternative, she believes its negation. We show that the testimonial is derived by the same two steps.

\section{The derivation}

Let $\phi$ be the letter and $\psi$ the sentence \enquote{Mr X is a good philosopher.} The purpose of the exchange, which both parties know, is to assess Mr X for a philosophy post; a letter of this kind is expected to address $\psi$. Hence $\phi \wedge \psi$, the letter the writer was expected to write, is an alternative of $\phi$ in the sense of Definition 3.4 of LNG-11-a: it is equally relevant, indeed more so, and stronger. It is more complex than $\phi$ by one clause, but the complexity clause of Definition 3.4 serves Manner, which forbids penalizing a speaker for not saying what was more cumbersome to say; it does not exclude a clause that the genre of the exchange requires.

\begin{enumerate}
\item The writer is presumed to observe the maxims at the level of what she conveys. The flout is open: it is an apparent violation of Quantity in what is said, which the reader resolves by assuming that the total signification is cooperative.
\item By Proposition 5.1 of LNG-11-a, she does not believe $\phi \wedge \psi$. Since she believes $\phi$, which she asserted, she does not believe $\psi$.
\item She is competent with respect to $\psi$: she has taught Mr X, and the reader knows it. This is the competence assumption, Assumption 3.5 of LNG-11-a.
\item By Proposition 5.2 of LNG-11-a, she believes not-$\psi$.
\end{enumerate}

No lexical scale is involved. Nothing in \enquote{excellent} or \enquote{regular} is a weaker member of an ordering whose stronger member concerns philosophical ability. The alternative $\psi$ is supplied by what the exchange is about, as \citet{Hirschberg1985} argued for the scales of scalar implicature generally, which she took to be partially ordered sets fixed by context rather than by the lexicon.

\section{Competence and cancellation}

The two steps come apart in a variant. Suppose the writer adds: \enquote{I have taught him only in the language classes, and have not seen his philosophical work.} The secondary implicature disappears, since the writer is not competent with respect to $\psi$; the primary implicature remains, since she still does not claim that he is a good philosopher. The reader learns nothing about Mr X's philosophy and does not conclude that he is bad at it. The added sentence cancels the secondary implicature without contradiction, as Theorem 4.2 of LNG-11-a requires of conversational implicatures.

\begin{objection}
A grammatical theory can generate the same reading: an exhaustivity operator applied to the letter, with alternatives that include contextually salient material \citep{Katzir2007}, yields \enquote{and nothing stronger relevant holds}, including that he is not a good philosopher.
\end{objection}

\begin{reply}
The grammatical operator applies whenever the alternative is in the set, and it negates it. In the variant, $\psi$ is as salient as before, and the letter's purpose is unchanged; yet the inference to not-$\psi$ is absent. What changed is what the writer knows. A grammatical theory can reproduce the variant only by making the alternative set depend on the speaker's competence, which is to import into the grammar the condition that Proposition 5.2 of LNG-11-a states. The case is thus a clear instance of the argument of LNG-11-a that the conditions under which these implicatures arise are conditions on speakers' knowledge.
\end{reply}

\section{Conclusion}

The testimonial is a Quantity implicature derived in the two Gricean steps. Its stronger alternative is supplied by the purpose of the exchange and not by the lexicon, and its secondary step depends on the writer's competence. Both features confirm that implicatures are calculable from the maxims and derived pragmatically.

\begin{thebibliography}{9}\small
\bibitem[Grice(1975)]{Grice1975} Grice, H.~P. (1975). Logic and conversation. In P. Cole and J.~L. Morgan (eds), \emph{Syntax and Semantics 3: Speech Acts}, 41--58. Academic Press.
\bibitem[Hirschberg(1985)]{Hirschberg1985} Hirschberg, J. (1985). \emph{A Theory of Scalar Implicature}. PhD dissertation, University of Pennsylvania.
\bibitem[Katzir(2007)]{Katzir2007} Katzir, R. (2007). Structurally-defined alternatives. \emph{Linguistics and Philosophy}, 30(6), 669--690.
\end{thebibliography}

\end{document}
